\begin{description}
\item[a)] The total angular momentum of the system is
\[ L = I\omega = (M_1 r_1^2 + M_2 r_2^2)\,\omega \]
\hfill(0.5 points)

We have also $M_1 r_1 = M_2 r_2$ and $D = r_1 + r_2$, which yield
\[ L = \frac{M_1 M_2}{M_1 + M_2} D^2 \omega \qquad (1) \]
\hfill(0.5 points)

The kinetic energy of the system is
\begin{align*}
K.E. &= \frac{1}{2} M_1 (r_1\omega)^2 + \frac{1}{2} M_2 (r_2\omega)^2
  = \frac{1}{2}\left(M_1 r_1^2 + M_2 r_2^2\right)\omega^2
  && \text{(0.5 points)} \\
&= \frac{1}{2}\,\frac{M_1 M_2}{(M_1 + M_2)}\, D^2\omega^2
  \qquad (2)
  && \text{(0.5 points)}
\end{align*}

\item[b)] From Newton's laws of motion we have
\[ M_1\omega^2 r_1 = M_2\omega^2 r_2 = \frac{G M_1 M_2}{D^2} \]
\hfill(1 point)

These equations together with those in a) yield
\[ \omega^2 = \frac{G(M_1 + M_2)}{D^3} \qquad (3) \]
\hfill(1 point)

\item[c)] In order to find the quantity $\Delta\omega$, we must also realize
that
\[ M_1 + M_2 = \text{constant} \qquad (4) \]
\hfill(0.5 points)

And that, since there is no external torque acting on the system, the total
angular momentum must be conserved.
\[ L = \frac{M_1 M_2}{M_1 + M_2} D^2\omega = \text{constant} \]
that is,
\[ M_1 M_2 D^2 \omega = \text{constant} \]
\hfill(0.5 points)

Now, after the mass transfer,
\begin{align*}
\omega &\to \omega + \Delta\omega \\
M_1 &\to M_1 + \Delta M_1 \\
M_2 &\to M_2 - \Delta M_1 \\
D &\to D + \Delta D
\end{align*}
Hence,
\[ M_1 M_2 D^2\omega
   = \left(M_1 + \Delta M_1\right)\left(M_2 - \Delta M_1\right)
     \left(D + \Delta D\right)^2\left(\omega + \Delta\omega\right) \]

After using the approximation $(1+x)^n \sim 1 + nx$ and rearranging, we get
\[ \frac{\Delta\omega}{\omega} + 2\,\frac{\Delta D}{D}
   = \left(\frac{M_1 - M_2}{M_1 M_2}\right)\Delta M_1 \qquad (5) \]
\hfill(1 point)

From equation (3), $\omega^2 D^3$ is also constant. That is,
\[ \omega^2 D^3 = (\omega + \Delta\omega)^2 (D + \Delta D)^3 \]
This gives,
\[ \frac{\Delta D}{D} = -\frac{2}{3}\,\frac{\Delta\omega}{\omega}
   \qquad (6) \]
\hfill(0.5 points)

Hence
\[ \Delta\omega = -\frac{3\left(M_1 - M_2\right)}{M_1 M_2}\,
   \omega\,\Delta M_1 \qquad (7) \]
\hfill(0.5 points)

\item[d)] Given that
\begin{gather*}
M_1 = 2.9\,\mathrm{M}_\odot\,, \quad M_2 = 1.4\,\mathrm{M}_\odot \\
\text{orbital period,}\quad T = 2.49\ \text{days}
\end{gather*}
and that $T$ has increased by 20\,s in the past 100 years, we have
\[ \omega = \frac{2\pi}{T} \quad\text{and}\quad
   \Delta\omega = -\frac{\omega}{T}\,\Delta T \qquad (8) \]
\hfill(0.5 points)

\[ \Delta M_1 = +\frac{1}{3}\left(\frac{M_1 M_2}{M_1 - M_2}\right)
   \frac{\Delta T}{T} \]
\hfill(0.5 points)
\begin{align*}
\frac{\Delta M_1}{M_1 \Delta t}
  &= \frac{1}{3}\left(\frac{1.4}{(2.9 - 1.4)}\right)
     \left(\frac{20}{2.49 \times 24 \times 3600}\right)
     \left(\frac{1}{100}\right) \\
  &= 2.89 \times 10^{-7}\ \text{per year}
  && \text{(0.5 points)}
\end{align*}

\item[e)] Mass is flowing from $M_2$ to $M_1$. \hfill(0.5 points)

\item[f)] From equations (6) and (8):
\[ \frac{\Delta D}{D\Delta t} = -\frac{2}{3}\,\frac{\Delta\omega}{\omega}
   = +\frac{2}{3}\,\frac{\Delta T}{T}
   = 6.20 \times 10^{-7}\ \text{per year} \]
\hfill(1.0 points)
\end{description}
